% Theory fragment for the parent REVTeX report. No standalone preamble.
% Derivations: theory/CLOSURE.md and theory/TRACE.md.
% Bibliography keys are supplied by the parent: MRD2018, BFSS2002, CLM2025.

\section{Original-equation certificates for the cyclic closure}
\label{sec:cyclic-theory}

The cyclic problem has indices modulo $n$ and original equations
\begin{align}
 G_i(x,q)&=q_i+\frac14q_{i+1}^{2}-x_i=0,\label{eq:cyclic-original}\\
 \dot x_i=f_i(x,q)&=-x_i+\frac{q_i+q_{i+1}}{16}+c_i,
 \qquad c_i=\frac1{32}+\frac{i\bmod3}{256}.
\end{align}
The declared branch starts at $x=q=0$, with physical time $0\le t\le1$,
and domains $X=[-1/8,1/8]^n$, $Q=[-1/4,1/4]^n$.
The native specialization admits $n=4,8,16$.
The required error is the supremum over the entire time interval and all
slow and algebraic coordinates, in the original infinity norm, with cap
$10^{-8}=1/100000000$.

\paragraph{Proposition C1: branch and physical-error transfer.}
There is a unique $q(x)\in Q$ for each $x\in X$. Moreover,
\begin{align}
 \|q(x)-q(z)\|_\infty&\le\frac87\|x-z\|_\infty,\label{eq:cyclic-rootlip}\\
 \|\widehat q-q(x)\|_\infty&\le\frac87\|G(x,\widehat q)\|_\infty
 \quad(\widehat q\in Q).
\end{align}
The reduced field $F(x)=f(x,q(x))$ contracts at rate $\lambda=6/7$.

\emph{Proof.}
The map $T_x(q)_i=x_i-q_{i+1}^2/4$ maps $Q$ into
$[-9/64,1/8]^n\subset Q$ and contracts with factor $1/8$.
Subtracting its fixed-point equations, or comparing $\widehat q$ to its
fixed point, gives the two bounds. On an upper face of $X$, $f_i\le-7/128$;
on a lower face, $f_i\ge1/8$. Thus the exact slow solution stays in $X$.
Writing $F=-x+Cq+c$ with $\|C\|_\infty=1/8$ gives the maximum-coordinate
Dini contraction rate $1-(1/8)(8/7)=6/7$.

For a differentiable candidate $(\widehat X,\widehat Q)$ in $X\times Q$,
define the original residuals
\begin{equation}
 R_f=\widehat X'-f(\widehat X,\widehat Q),\qquad
 R_g=G(\widehat X,\widehat Q).
\end{equation}
Whole-slab bounds $\|R_f\|_\infty\le D_f$ and
$\|R_g\|_\infty\le D_g$, with
$e=\|\widehat X-x_{\rm exact}\|_\infty$, imply
\begin{align}
 D^+e&\le-\lambda e+\delta,
 \qquad\delta=D_f+D_g/7,\label{eq:cyclic-error-ode}\\
 \|\widehat Q-q(x_{\rm exact})\|_\infty
 &\le\frac87(e+D_g).
\end{align}
This comparison uses a residual evaluated on the numerical reconstruction.
It therefore retains the exact-flow error anchor; an exact-input-only
discrete truncation estimate cannot be substituted without a separate
stability argument. Adjoint DAE analysis for linear quantities of interest
provides complementary output contracts~\cite{CLM2025}; the present contract
retains every physical coordinate throughout the path.

\subsection{Finite lift and the exact reconstructed path}

Let $b_i=q_{i+1}/2$, $(Av)_i=b_i v_{i+1}$ and $P=\prod_i b_i$.
Then $A^n=PI$ and
\begin{equation}
 (I+A)\sum_{k=0}^{n-1}(-A)^k=[1-(-1)^nP]I.
\end{equation}
The denominator $d=1-(-1)^nP$ satisfies $d\ge1-8^{-n}>0$ on $Q$.
Solving $v_i=f_i-b_i v_{i+1}$ by composition around the cycle and a backward
substitution computes the inverse action in $O(n)$ scalar operations.
For the exact rational lift $s(y)=(f,v)$, $y=(x,q)$,
\begin{equation}
 G_qv=f                                                        \label{eq:cyclic-tangency}
\end{equation}
holds even at endpoints with $G\ne0$.
The native circuit evaluates this declared analytic structure without a
Jacobian, JVP, AD or Newton callback. The identity is specific to these
equations and is not a black-box result for arbitrary closures.

Stored finite binary64 endpoints $y_0,y_1$ denote exact dyadics.
The certified Hermite polynomial uses the \emph{exact rational} slopes
$s_j=s(y_j)$, enclosed numerically by outward arithmetic. With
$B_j^m(\theta)=\binom mj\theta^j(1-\theta)^{m-j}$,
its cubic Bernstein controls are
\begin{equation}
 Y_0=y_0,\quad Y_1=y_0+hs_0/3,\quad
 Y_2=y_1-hs_1/3,\quad Y_3=y_1.                              \label{eq:cyclic-controls}
\end{equation}
The componentwise control hull must lie in $X\times Q$.
The Bernstein convex-hull property provides the whole-interval range
enclosure~\cite{MRD2018}. Arithmetic enclosure remains necessary: rounded
endpoint slopes would define a different path. The supplied dense-query
interface returns an interval enclosure of the declared path.

\paragraph{Proposition C2: sharp slow-residual weight.}
For one slow component let $X_j$ denote its controls and set
\begin{align}
 D_0&=f_i(y_0),\qquad D_2=f_i(y_1),\\
 D_1&=3(X_2-X_1)/h,\\
 r_1&=D_0/3+2D_1/3-f_i(Y_1),\\
 r_2&=2D_1/3+D_2/3-f_i(Y_2).
\end{align}
The cubic Bernstein residual coefficients are exactly $(0,r_1,r_2,0)$,
and consequently
\begin{equation}
 D_f\le\frac34\max_i\{|r_{1,i}|,|r_{2,i}|\}.               \label{eq:cyclic-df}
\end{equation}
\emph{Proof.}
Degree elevation of the quadratic derivative gives the displayed interior
coefficients. Affineness of $f$ gives field controls $f(Y_j)$; the endpoint
differences vanish by the exact-slope definition. The residual equals
$3\theta(1-\theta)[(1-\theta)r_1+\theta r_2]$.
Since $3\theta(1-\theta)\le3/4$, Eq.~\eqref{eq:cyclic-df} follows, with
equality possible for equal, same-signed $r_1,r_2$ at $\theta=1/2$.

\paragraph{Proposition C3: three interior algebraic coefficients.}
For component $i$ set $d_j=Q_{i,j}-X_{i,j}$, $z_j=Q_{i+1,j}$ and
$g_\ell=G_i(y_\ell)$. The degree-six residual coefficients are
\begin{equation}
 (g_0,g_0,\beta_2,\beta_3,\beta_4,g_1,g_1),\qquad
 \beta_k=E_k+S_k/4,                                      \label{eq:cyclic-gcoeff}
\end{equation}
where the exact rational expressions are
\begin{align}
 E_2&=(d_0+3d_1+d_2)/5,&S_2&=(2z_0z_2+3z_1^2)/5,\\
 E_3&=(d_0+9d_1+9d_2+d_3)/20,&
 S_3&=(z_0z_3+9z_1z_2)/10,\\
 E_4&=(d_1+3d_2+d_3)/5,&S_4&=(2z_1z_3+3z_2^2)/5.
\end{align}
\emph{Proof.}
Cubic-to-degree-six elevation gives $E_k$ and the Bernstein product rule
gives $S_k$. Equation~\eqref{eq:cyclic-tangency} makes the two endpoint
derivatives of $G(Y(\theta))$ zero. The derivative formulas
$6(\beta_1-\beta_0)$ and $6(\beta_6-\beta_5)$ then give the repeated
endpoint coefficients, without assuming exact constraint satisfaction.

Let $g$ bound the endpoint magnitudes and $m$ the three interior magnitudes,
over all components. The interior basis weight is
\begin{equation}
 w=B_2^6+B_3^6+B_4^6=15u^2-10u^3,
 \qquad u=\theta(1-\theta)\in[0,1/4].
\end{equation}
It increases with $u$ and attains $25/32$. Hence
\begin{equation}
 D_g\le\max\{g,(7g+25m)/32\}.                             \label{eq:cyclic-dg}
\end{equation}
The weight is sharp for this coefficient-information class. All these
constants follow from exact polynomial identities; none is fitted.

\paragraph{Proposition C4: endpoint reuse.}
The same represented endpoint and model/arithmetic identity define the
same exact slope and original residual. Thus $N$ consecutive slabs need
$N+1$ endpoint preparations instead of $2N$.
The native reuse check requires every endpoint bit, dimension and model
epoch to agree; branch, domain and arithmetic provider are fixed by the
module. For a nonautonomous extension the physical time would also belong
to this identity. A streaming consumer needs only the adjacent prepared
endpoints. This is an exact work reduction, with runtime effects measured
separately.

\subsection{A safe exponential-free error recurrence}

\paragraph{Proposition C5: maximum of two rational envelopes.}
Let $E,\delta,h\ge0$, $\lambda>0$ and $z=\lambda h$. Define
\begin{align}
 U_1&=\frac{E+h\delta}{1+z},&
 U_2&=\frac{(1-z/2)E+h\delta}{1+z/2},\\
 E_+&=\max(U_1,U_2),&E_{\rm whole}&=\max(E,E_+).            \label{eq:cyclic-rational-radius}
\end{align}
These enclose the endpoint and whole-slab comparison solution of
Eq.~\eqref{eq:cyclic-error-ode}.

\emph{Proof.}
Writing $w=\delta/\lambda$, the exact comparison is
$T(h)=w+(E-w)e^{-z}$. For $z\ge0$,
\begin{equation}
 \frac{1-z/2}{1+z/2}\le e^{-z}\le\frac1{1+z}.
\end{equation}
The upper inequality follows from $e^z\ge1+z$. For $0\le z<2$ the lower
one follows by integrating
$[\log((1+z/2)/(1-z/2))]'=1/(1-z^2/4)\ge1$; for $z\ge2$ its left side
is nonpositive. Use the upper exponential bound if $E\ge w$ and the lower
one if $E\le w$. The identity
\begin{equation}
 U_1-U_2=\frac{hz(\lambda E-\delta)}{2(1+z)(1+z/2)}
\end{equation}
shows that the maximum selects the safe expression. The scalar comparison
is monotone within the slab, yielding the whole-slab maximum.
In particular $U_1$ alone is invalid at $E=0$, $\delta>0$, $h>0$.

The independently valid inherited alternative is
$E_+=E/(1+\lambda h)+h\delta$, $E_{\rm whole}=E+h\delta$.
The two recipes are separate native comparison arms. The final algebraic
radius is $(8/7)(E_{\rm whole}+D_g)$. Rational constants, coefficients and
radius updates are evaluated outward under the core arithmetic assumptions;
nonfinite inputs and failed arithmetic are refused. The dyadic physical
clock has no uncharged time-label displacement. Admission compares with
the exact rational cap, rather than enlarging it by binary64 rounding.

\section{Typed quotient observables through a discriminant}
\label{sec:trace-theory}

\paragraph{Theorem T1: fixed rank with multiplicity.}
For a monic degree-$d$ polynomial $P$ and a polynomial observable $h$,
the multiplication operator on $\mathcal A=\mathbb C[z]/(P)$ satisfies
\begin{equation}
 \operatorname{Tr}_{\mathcal A}(m_h)
 =\sum_{r:P(r)=0}m_r h(r),\qquad \sum_r m_r=d.              \label{eq:trace-multiplicity}
\end{equation}
No square-free condition is needed. Monic division gives a fixed free basis
$1,z,\ldots,z^{d-1}$ even before parameter specialization. In the local
factor $\mathbb C[z]/((z-r)^{m_r})$, multiplication by $h$ is triangular
with diagonal $h(r)$ in the basis of powers of $z-r$. Its trace is
$m_rh(r)$; summing local factors proves the formula. If the coefficients of $P$ and $h$ depend polynomially on the parameters,
polynomial parameter dependence follows directly from monic reduction. This is the classical
trace/power-sum relation~\cite{BFSS2002}, specialized here with explicit
query semantics and finite-arithmetic obligations.

For the original two-equation family
$q_1^2+q_2=a$, $q_2^2+q_1=b$, elimination without division gives
\begin{equation}
 P_{a,b}(z)=z^4-2az^2+z+a^2-b,
 \qquad q_1=z,\quad q_2=a-z^2.                            \label{eq:trace-quartic}
\end{equation}
The quotient is therefore free of rank four at every $(a,b)$.
Its power sums are
\begin{align}
 &(S_0,S_1,S_2,S_3)=(4,0,4a,-3),\\
 &S_4=4a^2+4b,\quad S_5=-10a,\quad S_6=4a^3+12ab+3.
\end{align}
They follow from Newton identities, or directly by monic reduction and
Eq.~\eqref{eq:trace-multiplicity}. For $k>4$,
$S_k=2aS_{k-2}-S_{k-3}-(a^2-b)S_{k-4}$.
In particular $\operatorname{Tr}q_1=\operatorname{Tr}q_2=0$,
$\operatorname{Tr}(q_1q_2)=3$, and
$\operatorname{Tr}(q_1^2q_2^2)=4ab+3$.

At $a=b=3/4$,
\begin{equation}
 P=(z-1/2)^3(z+3/2),\qquad
 S_k=3(1/2)^k+(-3/2)^k.                                  \label{eq:trace-triple}
\end{equation}
The triple point contributes length three. Deduplicating roots changes
$S_0$ from $4$ to $2$ and $S_1$ from $0$ to $-1$.
The accepted query explicitly requests all complex roots with algebraic
multiplicity. A selected real branch is a different target: at $(a,b)=(2,2)$
the regular branch $(q_1,q_2)=(1,1)$ has determinant $4q_1q_2-1=3$ and
$q_1=1$, whereas both $\operatorname{Tr}q_1$ and its average are zero.
Thus a trace cannot be substituted as that branch's velocity, even away
from a singularity. The native interface also refuses real-only semantics.

\subsection{Shared field and a proved physical contraction constant}

The declared aggregate field is
\begin{align}
 \mathcal F_1&=-a+S_2/32+S_4/128+S_6/512,\\
 \mathcal F_2&=-b+S_2/64+S_4/64+S_6/256.
\end{align}
Exact substitution yields a single shared nonlinear circuit,
\begin{align}
 N(a,b)&=\frac{4a^3+16a^2+12ab+16b+3}{512},\\
 \mathcal F_1&=-7a/8+N,&\mathcal F_2&=-b+a/16+2N.          \label{eq:trace-shared}
\end{align}
Consequently $\mathcal F_2-2\mathcal F_1=29a/16-b$.
An evaluation error $e_N$ in the shared circuit contributes exactly
$(e_N,2e_N)$ before surrounding-operation errors. Outward interval
evaluation encloses the original polynomial for the declared parameter
intervals and preserves the same query. No inverse root separation enters.

\paragraph{Proposition T3: sharp unweighted row-bound constant.}
On $D=[1/2,1]^2$, the off-diagonal field derivatives are nonnegative and the
diagonals negative. Their row sums increase with $a,b$. At $(1,1)$,
\begin{equation}
 D\mathcal F=
 \begin{pmatrix}-49/64&7/128\\9/32&-57/64\end{pmatrix}.
\end{equation}
The row sums are $-91/128$ and $-39/64$. Therefore
\begin{equation}
 \mu_\infty(D\mathcal F)\le-39/64.                        \label{eq:trace-rate}
\end{equation}
The constant is sharp for this row-sum proof on this domain, since the
second row attains equality. Differentiation supplies an offline proof of
the bound; the runtime field uses only the polynomial circuit.

The controlled original target used to cross the discriminant is
\begin{align}
 p(t)&=\frac34+\frac t8+\frac{t^5}{64},\qquad -1\le t\le1,\\
 \dot y&=\mathcal F(y)+p'(t)\mathbf1-\mathcal F(p(t)\mathbf1),\\
 y(-1)&=(39/64)\mathbf1.                                  \label{eq:trace-controlled}
\end{align}
Its exact solution is $y=p(t)\mathbf1$, with terminal value
$(57/64)\mathbf1$. Since $p'>0$ and
$\operatorname{disc}P_{a,a}=(4a-3)^3(4a+1)$, it crosses the triple-root
parameter at $t=0$. The time-only forcing preserves
Eq.~\eqref{eq:trace-rate}. A candidate's uniform original-field defect
$\delta$ and a certified domain tube then give
\begin{equation}
 e(t_i+s)\le e^{-(39/64)s}E+
 \frac{1-e^{-(39/64)s}}{39/64}\delta .                     \label{eq:trace-physical}
\end{equation}
A strict boundary margin with a whole-slab error bound closes the usual
first-exit argument when domain containment is not already known.
Trace arithmetic included in an original-field interval residual is charged
once. Endpoint agreement with $p(t)$ is an independent control and does not
replace the whole-path residual/tube obligation.

\section{Boundaries of the admitted mathematical scope}
\label{sec:theory-boundaries}

The selected-branch atlas and the multiplicity-preserving trace have
different mathematical targets. The two supplied primitive charts share
the determinant $4q_1q_2-1$; their transition formulas do not establish
regular selected-root continuation at determinant zero. The cyclic domain
instead has the strictly positive denominator proved above. Trace regularity
at Eq.~\eqref{eq:trace-triple} follows from its monic rank-four algebra and
does not resolve a selected singular branch.

The separate F5 global incidence counts also remain open. The original
$116/58$ ordered/unordered totals omit a horizontal contribution $4/2$.
The resulting $120/60$ remains conditional on completeness of the retained
charts and exclusions. A reviewed transversality stratum with
$D=4\eta^3p+27\eta^2q+27pq=0$ has local length three despite regular
carrier data, so reduced point counting alone cannot certify the quotient
rank. A future multivariate trace provider must supply its full ideal,
localizations, basis or finite-freeness proof, local lengths and quotient
semantics. These incomplete global counts do not affect the independent
rank-four proof for Eq.~\eqref{eq:trace-quartic}.

The native components instantiate the bounded contracts proved here;
their measured validation is reported separately. Whole-model physical
certification for a different application still requires its original
closure/ODE, selected observable, domain, physical clock and representation
errors. Generic stiff integration, singular-branch continuation, arbitrary
sparse quotients and wider floating-point environments retain independent
theoretical and implementation obligations. No generic production claim
follows solely from the fixed-model certificates.
